\addchap{Abstract}
\noindent A law thesis needs an identifiable research question, a defensible
method and a contribution that the evidence can support. This working example
shows how to organise those elements in a single document. It demonstrates
chapter-based footnotes, repeat references, jurisdiction-specific tables of
authorities and a secondary-source bibliography. Its chapters provide prompts
for doctrinal, comparative and empirical inquiry, together with an alternative
structure for a thesis incorporating published articles. The accompanying
records distinguish source retrieval, legal interpretation and verification
of particular propositions. Separate layouts support reading, supervisor
review and printing. A Groningen profile records institution-specific matters
for candidates to confirm. The example supplies a usable starting point for
writing and review; the candidate remains responsible for the argument,
citations, permissions and applicable submission requirements.

\LawPrompt{Replace this example with your question, method, principal findings
and original contribution. Check your university's abstract word limit and
language requirements. Write out abbreviations; omit footnotes here.}
